\section{A Cautionary Note on Trust Regions}
\label{sec:trust}

The retired configuration projected the student's complete parameter vector
back into a relative-$L^2$ ball around the released M7 parameters after every
optimizer update:
\begin{equation}
\theta \leftarrow \theta_0 +
\min\!\left(1,\frac{r\lVert\theta_0\rVert_2}
{\lVert\theta-\theta_0\rVert_2}\right)(\theta-\theta_0),
\qquad r=\ballRadius.
\label{eq:trust}
\end{equation}
We describe it separately from the retired loss terms of
Appendix~\ref{sec:also-tried} because its consequences were not local to itself.
While it was active, it made the other terms unmeasurable, and every ablation
verdict recorded under it is uninformative.

\subsection{Why It Seemed Reasonable}

The premise was conservative and, stated abstractly, sound.
The released M7 is a working expert; the fine teacher is locally more
informative but globally incomplete and model-derived; a student free to move
anywhere in parameter space can follow a confident teacher into a region where
it is wrong, and the loss will not notice because the teacher supplies the
targets.
Bounding the student to a neighbourhood of a known-good function is a cheap
insurance policy against that failure.
The radius was not arbitrary either: it was derived from an earlier measurement
of how far the function could safely depart, and inherited across several
subsequent versions.

The mistake was not the mechanism.
It was retaining a radius across a change of objective, corpus, and optimizer
without re-measuring what the new configuration actually required, and then
reading the resulting absence of movement as an absence of effect.

\subsection{What It Actually Did}

The telemetry was unambiguous and was recorded at the time.
The projection was active for $\ballActiveFirst$\% of updates in the first
evaluation interval and for 100\% of updates in every interval thereafter, with
a mean pre-projection step slightly larger than the radius itself.
Roughly a quarter of every optimizer step was being discarded.
The ball was not a guardrail that occasionally caught an excursion; it was a
boundary the run rode continuously, which makes the radius a hyperparameter of
the method rather than a safety margin.

The decisive number came later, and from an observation rather than an
experiment.
Once a wide-corpus run was trained without any projection, we simply measured
where it ended up.
The recipe that produced the model in this paper sits at roughly
$\ballRatioThirty$ times the shipped radius after $30{,}000$ exposures and
$\ballRatioFifty$ times after $50{,}000$, and it keeps moving.
The displacement the successful method requires is nearly two orders of
magnitude outside the ball that earlier versions were confined to.
No amount of loss design inside that ball could have found this model, because
the model is not in it.

\subsection{The Second-Order Damage}

A saturated constraint does not merely slow a run down.
It silently converts every comparison made under it into a comparison of
projections.
Two arms differing in one loss term will produce nearly identical functions if
the projection removes most of the difference, and the correct reading of
``these two runs are indistinguishable'' is then ``this experiment had no
power,'' not ``this term does nothing.''
Our own record contains the clean example: a carefully matched pair testing one
mask refinement of the M7 anchor produced probability fields differing by
fractions of a percent at every threshold, and the run logs show why---the
anchor's own loss component was mechanically near zero, because the function
never left M7's neighbourhood.
That pair was recorded as a null result for the refinement.
The refinement is in the shipped objective today, and it is one of the three
mask properties in Sec.~\ref{sec:crossres} that the term needs to work.

This is the mechanism by which a saturated trust region manufactures a stack of
loss terms.
Each new term is added, measured as harmless, and kept; the terms that would
have been rejected are protected by the same constraint that prevents the terms
that would have been accepted from showing a benefit.

\subsection{What We Would Do Instead}

We do not conclude that parameter-space trust regions are a bad idea.
We conclude that an unmeasured one is, and that the measurement is cheap.

Log the projection's activity fraction and the pre-projection step norm from the
first interval, and treat sustained saturation as a failed run rather than a
successful guard.
Before trusting any ablation, measure the unconstrained displacement of the same
recipe and compare it with the radius; if the free run ends up far outside, the
ablation was run at zero power and its verdict should be discarded rather than
recorded.
Re-derive the radius whenever the objective, corpus, or optimizer changes, since
a radius calibrated for one of them carries no authority over another.
Where the underlying worry is following a model-derived teacher into its own
errors, prefer a constraint expressed in function space and restricted to where
the teacher is silent or confidently agrees, which is what
Eq.~\ref{eq:m7-kl} does, and which survived ablation when the parameter-space
constraint did not.
